\chapter*{Abstract}
\addcontentsline{toc}{chapter}{Abstract}
\thispagestyle{plain}

Electrochemical detection methods, particularly amperometry and chronoamperometry, offer exceptional sensitivity, selectivity, and miniaturisation potential for point-of-care biosensing. However, laboratory-grade potentiostats remain prohibitively expensive---commercial instruments such as the PalmSense EmStat series cost approximately \texttildelow\rupee{}5,00,000---placing them out of reach for resource-limited research settings. This thesis presents the design, fabrication, firmware development, and multi-analyte validation of a low-cost, portable, web-connected amperometric device aimed at democratising electrochemical sensing.

The hardware platform is built around the ESP32-WROOM-32E System-on-Chip, which provides a dual-core 240\,MHz processor together with integrated 802.11\,b/g/n Wi-Fi connectivity. A 16-bit dual-channel digital-to-analogue converter (AD5667R, Analog Devices) provides precise potential control, while a quad-channel precision operational amplifier (AD8608, Analog Devices) implements a transimpedance amplifier (TIA) for current-to-voltage conversion and a voltage follower for reference electrode buffering. Four selectable gain resistors switched via MAX4643/MAX4644 analogue switches enable current measurement across multiple orders of magnitude. Power management is handled by a MCP73831 lithium-polymer charger and dual MIC5205 3.3\,V low-dropout regulators with a star-point grounding topology to minimise analogue-digital interference. The printed circuit board was designed in EasyEDA and fabricated by JLCPCB, yielding a total bill-of-materials cost of approximately \texttildelow\rupee{}5,000--7,000.

The firmware, developed using the ESP32 Arduino framework, implements a complete chronoamperometry routine with user-configurable applied potential, acquisition duration, sampling interval, and current range. A single-page web application served directly from the microcontroller enables wireless real-time visualisation via Chart.js, with export options in CSV and JSON formats. An enclosure was designed in Autodesk Fusion\,360 and printed in PLA using a Prusa Mini FDM printer.

Electrochemical performance was validated using four analyte-electrode systems: (i)~ferri/ferrocyanide \ferrocene{} as an outer-sphere redox probe for basic device characterisation; (ii)~ruthenium hexamine \ruhex{} as a surface-insensitive probe; (iii)~hydrogen peroxide (\ce{H2O2}) at gold screen-printed electrodes; and (iv)~p-cresol, a clinically relevant uremic toxin, over a concentration range spanning \SIrange{1e-15}{1e-2}{\molar}. Cottrell analysis of chronoamperometric transients confirmed the device's ability to resolve diffusion-controlled faradaic currents consistent with literature-reported values. The system achieves sub-femtomolar detection of p-cresol, demonstrating ultra-high sensitivity competitive with laboratory instruments at a fraction of the cost.

\vspace{0.4cm}
\noindent\textbf{Keywords:} amperometric biosensor, chronoamperometry, ESP32, low-cost potentiostat, web-based monitoring, p-cresol, hydrogen peroxide, point-of-care diagnostics.
